%!TEX root = ../main.tex

\chapter{Análisis de sostenibilidad e implicaciones}
\label{ch:impacto}

\section{Impacto ambiental}
\label{sec:impacto_ambiental}

La mayor parte del impacto ambiental del sistema proviene de los componentes electrónicos, los módulos comerciales y la placa de circuito impreso. Al ser un equipo único, su huella de fabricación es pequeña, pero le afectan los mismos factores que a cualquier componente electrónico. Por un lado, los materiales usados están sujetos a la normativa RoHS \cite{dir_rohs}, la cual limita el uso de sustancias como el plomo o el mercurio. Por otro lado, al final de su vida útil, el conjunto debe ser gestionado como residuo de aparato eléctrico y electrónico (RAEE) \cite{dir_raee}, requiriendo un reciclaje separado de la basura común.

En el ámbito de su alimentación y consumo, el gasto del sistema en funcionamiento es muy bajo, ofreciendo un impacto mínimo en este aspecto. Al alimentarse por USB-C y no incluir baterías, tampoco genera residuos asociados con las celdas de litio u otros materiales presentes en ellas.

Por último, el hecho de que el software sea de código abierto ayuda a alargar la vida del producto, ya que cualquiera puede repararlo, modificarlo o reutilizarlo, reduciendo así la obsolescencia y evitando su desecho antes de tiempo.